% !TEX program = xelatex
\documentclass[11pt,a4paper]{article}

\usepackage[margin=2cm]{geometry}
\usepackage{fontspec}
\usepackage{xcolor}
\usepackage{tikz}
\usepackage{enumitem}
\usepackage{tabularx}
\usepackage{hyperref}
\usepackage{microtype}
\usepackage{setspace}
\usepackage{titlesec}
\usepackage{fancyhdr}
\usepackage{parskip}
\usepackage{graphicx}
\usepackage{calc}
\usepackage{array}

% —— Stone Soft ——
\definecolor{ink}{HTML}{1C1917}
\definecolor{mute}{HTML}{78716C}
\definecolor{soft}{HTML}{A8A29E}
\definecolor{line}{HTML}{D6D3D1}
\definecolor{paper}{HTML}{FAFAF9}
\definecolor{subtitle}{HTML}{4A4A4A}
\definecolor{graphite}{HTML}{0C0A09}
\definecolor{wash}{HTML}{F5F5F4}

\setmainfont{TeX Gyre Pagella}
\setsansfont{Montserrat}[
  UprightFont = * ,
  BoldFont = * Medium ,
  LetterSpace = 2.0
]
\newfontfamily\fDisplay{TeX Gyre Pagella}
\newfontfamily\fSans{Montserrat}
\newfontfamily\fSansMed{Montserrat Medium}

\hypersetup{
  colorlinks=true,
  linkcolor=ink,
  urlcolor=mute,
  pdfauthor={Chiiko · Sophie Gaëlle Gomez},
  pdftitle={Press Kit — Sophie Gaëlle Gomez (ejemplo Stone Soft)},
  pdfsubject={Ejemplo de organización de press kit}
}

\pagecolor{paper}
\color{ink}
\pagestyle{fancy}
\fancyhf{}
\renewcommand{\headrulewidth}{0pt}
\renewcommand{\footrulewidth}{0.4pt}
\renewcommand{\footrule}{\hbox to\headwidth{\color{line}\leaders\hrule height \footrulewidth\hfill}}
\fancyfoot[L]{\footnotesize\color{mute}\fSans Ejemplo · Press kit · Stone Soft}
\fancyfoot[R]{\footnotesize\color{mute}\thepage}

\titleformat{\section}
  {\fDisplay\color{ink}\LARGE}
  {}{0em}{}
  [\vspace{0.15em}{\color{line}\titlerule[0.6pt]}]
\titleformat{\subsection}
  {\fSansMed\color{subtitle}\footnotesize}
  {}{0em}{\MakeUppercase}

\newcommand{\hint}[1]{%
  \par\vspace{0.35em}%
  \noindent{\footnotesize\color{mute}\fSans\textit{→ Qué irá aquí:}~#1}\par%
  \vspace{0.5em}%
}

\newcommand{\sectionlabel}[1]{%
  {\fSansMed\color{soft}\footnotesize\MakeUppercase{#1}}\par\vspace{0.4em}%
}

\newcommand{\photobox}[1]{%
  \begin{tikzpicture}
    \fill[wash] (0,0) rectangle (7.2,4.6);
    \draw[line, line width=0.6pt] (0,0) rectangle (7.2,4.6);
    \node[align=center, text=mute, font=\fSans\footnotesize] at (3.6,2.3) {#1};
  \end{tikzpicture}%
}

\newcommand{\swatch}[2]{%
  \begin{tikzpicture}[baseline=(s.base)]
    \fill[#1] (0,0) rectangle (0.85,0.85);
    \draw[line, line width=0.4pt] (0,0) rectangle (0.85,0.85);
    \node[anchor=west, font=\fSans\tiny\color{mute}, inner sep=0] (s) at (1.0,0.42) {#2};
  \end{tikzpicture}%
}

\begin{document}
\thispagestyle{empty}

% ===================== COVER =====================
\begin{tikzpicture}[remember picture,overlay]
  \fill[paper] (current page.south west) rectangle (current page.north east);
  \fill[wash] ([yshift=-0cm]current page.north west) rectangle ([yshift=-4.2cm]current page.north east);
  \draw[line, line width=0.6pt]
    ([yshift=-4.2cm]current page.north west) -- ([yshift=-4.2cm]current page.north east);
\end{tikzpicture}

\vspace*{1.6cm}
{\fSansMed\color{soft}\footnotesize\MakeUppercase{Press kit · Ejemplo de organización}}

\vspace{1.1cm}
{\fDisplay\color{graphite}\fontsize{36}{40}\selectfont Sophie Gaëlle Gomez}

\vspace{0.55cm}
{\fSansMed\color{subtitle}\normalsize\MakeUppercase{Actrice · Modèle · Silver Presence}}

\vspace{0.9cm}
{\color{line}\rule{\textwidth}{0.6pt}}

\vspace{0.75cm}
\begin{minipage}[t]{0.55\textwidth}
  {\fSans\color{mute}\small
  Documento de muestra para definir la estructura del press kit digital y PDF.\\[0.4em]
  Paleta: \textbf{\color{ink}Stone Soft}\\
  Hex: \#1C1917 · \#78716C · \#A8A29E · \#FAFAF9
  }
\end{minipage}
\hfill
\begin{minipage}[t]{0.38\textwidth}
  \raggedleft
  {\fSans\color{mute}\footnotesize
  México · Europe\\
  sophiegomez.me\\[0.6em]
  \textbf{\color{ink}Versión ejemplo}\\
  textos guía + placeholders
  }
\end{minipage}

\vspace{1.6cm}
\sectionlabel{Paleta Stone Soft}
\vspace{0.3em}
\swatch{ink}{Ink \#1C1917}\hspace{0.8em}
\swatch{mute}{Mute \#78716C}\hspace{0.8em}
\swatch{soft}{Soft \#A8A29E}\hspace{0.8em}
\swatch{line}{Line \#D6D3D1}\hspace{0.8em}
\swatch{paper}{Paper \#FAFAF9}

\vspace{2.2cm}
{\fSans\color{mute}\footnotesize
Este PDF muestra \textbf{\color{ink}qué secciones} debe tener el press kit y, bajo cada bloque,
\textbf{\color{ink}qué contenido real} sustituirá a los textos de ejemplo.
}

\vfill
{\fSans\color{soft}\footnotesize Chiiko · Propuesta de organización · No indexar}

\newpage

% ===================== ORGANIZACIÓN =====================
\section*{Cómo está organizado}

\hint{Una página de índice interno para periodistas, casting y prensa. En la versión final puede omitirse o reducirse a un solo párrafo.}

\vspace{0.4em}
{\fSansMed\color{subtitle}\footnotesize\MakeUppercase{Mapa del press kit}}

\vspace{0.6em}
\renewcommand{\arraystretch}{1.45}
\begin{tabularx}{\textwidth}{@{} >{\fSans\footnotesize\bfseries\color{ink}}p{3.2cm} >{\small\color{ink}}X @{}}
\hline
\color{line}\rule{0pt}{2.4em}Portada & Nombre, roles, ciudad/base, sitio web y sello visual Stone Soft.\\
Ficha rápida & Datos esenciales en 8--12 líneas (idiomas, bases, estatura, contactos).\\
Bio corta & 80--120 palabras para notas de prensa y pie de foto.\\
Bio larga & 180--280 palabras para dossiers, festivales y medios.\\
Perfil técnico & Medidas / estatura / idiomas / formación (actriz + modelo).\\
Créditos & Filmografía y trabajos selectos (cine, TV, teatro, campaña).\\
Formación & Escuelas, coaches, workshops relevantes.\\
Silver Presence & 1 párrafo sobre el proyecto de voz / comunicación.\\
Fotos & Guía de headshots / stills en HD + créditos de fotógrafo.\\
Contacto & Agencia, email directo, Instagram, IMDb, web.\\
\hline
\end{tabularx}

\vspace{1.2em}
{\fSans\color{mute}\footnotesize
Orden recomendado para PDF descargable (4--8 páginas). En web, las mismas secciones pueden vivir como bloques en \texttt{/press-kit}.
}

\newpage

% ===================== FICHA RÁPIDA =====================
\section*{Ficha rápida}

\hint{Tarjeta de identidad profesional. Actualizar cuando cambien representación, bases o medidas.}

\vspace{0.5em}
\begin{tabularx}{\textwidth}{@{} >{\fSans\footnotesize\color{mute}}p{3.4cm} >{\small\color{ink}}X @{}}
Nombre & Sophie Gaëlle Gomez\\
Roles & Actriz · Modelo · Silver Presence\\
Bases & Ciudad de México · Europa (ciudades a confirmar)\\
Idiomas & Francés · Español · Inglés (nivel real a confirmar)\\
Estatura & [ej.\ 1,XX m] --- \textit{\color{mute}dato oficial de Sophie}\\
Agencia (MX) & Talent Boutique --- \textit{\color{mute}o la representación vigente}\\
Email & contact@sophiegomez.me --- \textit{\color{mute}email de bookings / prensa}\\
Web & \url{https://sophiegomez.me}\\
IMDb & \url{https://www.imdb.com/es/name/nm1960841/}\\
Instagram & @sophiegaelleof --- \textit{\color{mute}handle oficial}\\
\end{tabularx}

\newpage

% ===================== BIO CORTA =====================
\section*{Biografía corta}

\hint{Versión de 80--120 palabras. Sirve para press release, programas de mano, pies de foto y emails a casting. Tono sobrio, en tercera persona.}

\vspace{0.35em}
\sectionlabel{Texto de ejemplo (sustituir)}

{\large\color{ink}
Sophie Gaëlle Gomez es actriz y modelo entre México y Europa. Formada en escena clásica y contemporánea, construye una presencia sobria —silencio, cuerpo, precisión— en cine de autor, teatro de cámara e imagen editorial. Con \emph{Silver Presence}, desarrolla también un espacio de voz y relato. Trabaja en francés, español e inglés.
}

\vspace{0.9em}
{\fSans\color{mute}\footnotesize
\textbf{\color{ink}Checklist del texto final:} ciudad(es) base · 1--2 géneros que la definen · idiomas · mención breve de Silver Presence (opcional) · sin jerga promocional.
}

\newpage

% ===================== BIO LARGA =====================
\section*{Biografía larga}

\hint{Versión dossier (180--280 palabras). Para festivales, prensa cultural, partners y press kit completo. Puede existir en FR / ES / EN.}

\vspace{0.35em}
\sectionlabel{Texto de ejemplo (sustituir)}

{\color{ink}
Sophie Gaëlle Gomez construye una trayectoria entre la escena, el cine y la imagen. Como un lanzamiento: precisión, silencio, impacto. Entre México y Europa, habita cada cuadro —actriz, modelo, voz.

Formada en trabajo escénico clásico y contemporáneo, busca roles donde el cuerpo y la presencia importan tanto como la palabra. En imagen, trata la luz, la línea y la postura como lenguaje visual: editorial, belleza y campaña, con la misma exigencia estética.

\emph{Silver Presence} es su proyecto de comunicación: un espacio dedicado a la voz, la presencia y el relato, donde explora cómo se construyen la imagen y la palabra.

\textit{\color{mute}[Aquí añadir 2--4 créditos reales destacados en una frase, p.\ ej.\ «Recientemente participó en\ldots» o «Ha colaborado con\ldots».]}
}

\vspace{0.8em}
{\fSans\color{mute}\footnotesize
\textbf{\color{ink}Qué completar con Sophie:} formación concreta · créditos verificables · tono (más cine / más fashion) · versión en francés e inglés.
}

\newpage

% ===================== PERFIL TÉCNICO =====================
\section*{Perfil técnico}

\hint{Bloque útil para casting y bookings. Separar actriz y modelo si las medidas o usos son distintos.}

\vspace{0.4em}
\subsection*{Actriz}
\begin{tabularx}{\textwidth}{@{} >{\fSans\footnotesize\color{mute}}p{3.4cm} >{\small\color{ink}}X @{}}
Estatura & [1,XX m]\\
Tipo de rol & [dramático / contemporáneo / autor --- a definir]\\
Idiomas de juego & Francés · Español · Inglés\\
Bases & México · Europa\\
Demo reel & Enlace o QR al reel en \texttt{/actrice}\\
\end{tabularx}

\vspace{0.8em}
\subsection*{Modelo}
\begin{tabularx}{\textwidth}{@{} >{\fSans\footnotesize\color{mute}}p{3.4cm} >{\small\color{ink}}X @{}}
Estatura & [dato oficial]\\
Medidas & [busto / cintura / cadera --- solo si Sophie autoriza]\\
Cabello / ojos & [color real]\\
Book & Editorial · Comercial (enlaces al sitio)\\
\end{tabularx}

\vspace{0.9em}
{\fSans\color{mute}\footnotesize
\textbf{\color{ink}Importante:} no inventar medidas. Dejar celdas vacías hasta confirmación. El PDF final solo publica lo autorizado.
}

\newpage

% ===================== CRÉDITOS =====================
\section*{Créditos selectos}

\hint{Lista curada (no exhaustiva). Formato: Título — Rol — Director / Medio — Año. Priorizar calidad y verificabilidad (IMDb).}

\vspace{0.35em}
\sectionlabel{Ejemplo de formato}

\begin{itemize}[leftmargin=1.1em, itemsep=0.45em]
  \item {\color{ink}\textbf{Título del proyecto} — Rol (ej.\ «Ana») — Dir.\ Nombre Apellido — Cine / 20XX}\\
  {\footnotesize\color{mute}\textit{→ Sustituir por crédito real; 1 línea por obra.}}
  \item {\color{ink}\textbf{Obra de teatro} — Personaje — Compañía / Sala — 20XX}\\
  {\footnotesize\color{mute}\textit{→ Incluir solo montajes que Sophie quiera mostrar.}}
  \item {\color{ink}\textbf{Campaña / editorial} — Marca o revista — Fotógrafo — 20XX}\\
  {\footnotesize\color{mute}\textit{→ Para el book modelo: 4--8 trabajos representativos.}}
\end{itemize}

\vspace{0.7em}
{\fSans\color{mute}\footnotesize
Recomendación: máximo 8--12 créditos en el PDF. El resto vive en IMDb y en el sitio.
}

\newpage

% ===================== FORMACIÓN =====================
\section*{Formación}

\hint{Escuelas, estudios, coaches y workshops con nombre real. Evitar listas interminables: 4--8 entradas fuertes.}

\begin{itemize}[leftmargin=1.1em, itemsep=0.45em]
  \item {\color{ink}\textbf{[Escuela / estudio]} — Ciudad — años o programa}\\
  {\footnotesize\color{mute}\textit{→ Ej.\ formación teatral, cine, voz.}}
  \item {\color{ink}\textbf{[Coach / workshop]} — Especialidad — año}\\
  {\footnotesize\color{mute}\textit{→ Solo si aporta credibilidad internacional o técnica.}}
\end{itemize}

\newpage

% ===================== SILVER PRESENCE =====================
\section*{Silver Presence}

\hint{Un párrafo corto para medios y partners. No hace falta dossier aparte salvo que el proyecto crezca.}

\vspace{0.35em}
\sectionlabel{Texto de ejemplo (sustituir)}

{\color{ink}
Silver Presence es el proyecto de comunicación de Sophie Gaëlle Gomez: un espacio dedicado a la voz, la presencia y el relato. A través de video y contenido editorial, explora cómo se construyen la imagen y la palabra entre México y Europa.
}

\vspace{0.6em}
{\fSans\color{mute}\footnotesize
\textbf{\color{ink}Añadir en final:} enlace a \texttt{/silver-presence}, Instagram del proyecto (si existe), 1 frase de posicionamiento.
}

\newpage

% ===================== FOTOS =====================
\section*{Fotos de prensa}

\hint{En el PDF final: 4--6 fotos HD con crédito. Aquí solo se muestra la organización visual (placeholders).}

\vspace{0.5em}
\begin{minipage}[t]{0.48\textwidth}
  \photobox{Headshot 1\\[0.3em]Retrato frontal · luz natural\\[0.2em]{\tiny Crédito fotógrafo}}
  \vspace{0.35em}
  {\fSans\tiny\color{mute}→ Headshot principal para casting / prensa}
\end{minipage}
\hfill
\begin{minipage}[t]{0.48\textwidth}
  \photobox{Headshot 2\\[0.3em]3/4 o perfil · variación\\[0.2em]{\tiny Crédito fotógrafo}}
  \vspace{0.35em}
  {\fSans\tiny\color{mute}→ Variante de expresión / look}
\end{minipage}

\vspace{1.0em}
\begin{minipage}[t]{0.48\textwidth}
  \photobox{Still / escena\\[0.3em]Frame de proyecto\\[0.2em]{\tiny Autorización de uso}}
  \vspace{0.35em}
  {\fSans\tiny\color{mute}→ Still con derechos despejados}
\end{minipage}
\hfill
\begin{minipage}[t]{0.48\textwidth}
  \photobox{Editorial\\[0.3em]Look book / campaña\\[0.2em]{\tiny Crédito + marca}}
  \vspace{0.35em}
  {\fSans\tiny\color{mute}→ Imagen modelo (si aplica)}
\end{minipage}

\vspace{1.1em}
{\fSans\color{mute}\footnotesize
\textbf{\color{ink}Normas:} JPG/PNG alta resolución · espacio de color sRGB · incluir crédito · indicar si el uso es solo prensa o también comercial.
}

\newpage

% ===================== CONTACTO =====================
\section*{Contacto}

\hint{Cerrar el PDF con una sola página de contacto clara. Evitar demasiados emails.}

\vspace{0.5em}
\begin{tabularx}{\textwidth}{@{} >{\fSans\footnotesize\color{mute}}p{3.6cm} >{\small\color{ink}}X @{}}
Prensa / general & contact@sophiegomez.me\\
Acting — México & Talent Boutique · info@talentboutique.com\\
Instagram & \url{https://www.instagram.com/sophiegaelleof/}\\
IMDb & \url{https://www.imdb.com/es/name/nm1960841/}\\
Sitio & \url{https://sophiegomez.me}\\
Press kit web & \url{https://sophiegomez.me/press-kit}\\
\end{tabularx}

\vspace{1.4em}
{\color{line}\rule{\textwidth}{0.6pt}}

\vspace{0.9em}
{\fSansMed\color{soft}\footnotesize\MakeUppercase{Próximo paso con Sophie}}

\vspace{0.45em}
{\small\color{ink}
Reemplazar todos los bloques marcados con \textit{\color{mute}→ Qué irá aquí} por datos reales, elegir 4--6 fotos autorizadas, cerrar bio corta + larga en FR/ES/EN, y exportar la versión definitiva del PDF para \texttt{public/press-kit.pdf}.
}

\vspace{1.6em}
{\fSans\color{mute}\footnotesize
Documento de ejemplo · Paleta Stone Soft · Chiiko\\
No es el press kit final — es la plantilla de organización.
}

\end{document}
